% Section 9.1, from lectures/L09/README.md: the objectives, the evaluation questions, and the next
% lecture.

\needspace{10\baselineskip}
\section{Review}
\label{c9:sec:review}

\subsection*{What you should be able to do now}
After this chapter, you should:
\begin{itemize}
  \item Have created a concrete subclass \code{MaxPool} that inherits
    \code{ml::conv_layer::Interface}.
  \item Have implemented \code{feedforward()}, \code{backpropagate()}, and \code{optimize()} for
    \code{MaxPool}.
\end{itemize}

\needspace{6\baselineskip}
\subsection*{Questions to test yourself}
You should be able to explain:
\begin{enumerate}
  \item How does your \code{ml::conv_layer::MaxPool} implementation know which positions the
    gradients should be propagated back to during backpropagation?
  \item What differences exist between a max pooling layer and a conv or dense layer, in terms of
    trainable parameters?
  \item Why can \code{Conv} and \code{MaxPool} share a single interface
    (\code{conv_layer::Interface}) despite one having trainable parameters and the other not?
\end{enumerate}

\needspace{6\baselineskip}
\subsection*{Looking ahead}
Chapter~\ref{c10:ch:flatten} implements the flatten layer, and wires the complete CNN together as
the book's capstone.
